\section{与同期工作的关系}\label{app:comparison}

Qiushi Engine 独立形成了本文的结构性证明。Wang 的同期预印本建立了相同的数值下界\cite{wang2026twentyone}。我们在研究及原始科学报告完成之后的发布准备阶段，才获悉该工作。

下面对照 Wang 已公开的二元域证明与复现包，以及本文记录的证明与研究历程，区分共同的数学基础与有限前提、推导和研究贡献上的差异。

\paragraph{共同的数学基础。}
两种证明都继承了经过认证的限制下界，并使用占据容量；都导出了高秩第一因子两两之差秩一，并将它们置于共同的仿射行陪集或列陪集中。按零化子限制输入，与对系数空间取商，是同一运算在对偶约定下的表述。D'Ambrosio 的工作更早使用了容量论证与紧限制\cite{dambrosio2026}。

\begin{longtable}{@{}L{2.1cm}L{5.7cm}L{6.4cm}@{}}
\toprule
方面&Wang 的二元域证明&Qiushi Engine 的结构性证明\\
\midrule
\endfirsthead
\toprule 方面&Wang 的二元域证明&Qiushi Engine 的结构性证明\\\midrule
\endhead
\bottomrule\endfoot
有限前提
&通过秩一张成搜索与回溯加强限制表。
&以固定的早期表为基础，新增八个商下界，由 CNF/DRAT 反驳证明与闭合的 Lean 整数分支证书分别认证。\\
陪集归约之后
&分类 35 种高秩构型，并依据加强后的表检验其秩一补全。
&将仿射平面容量与完整张量的秩和结合，强迫得到唯一秩型 $(16,1,3)$，并在秩和 27 处取等。\\
决定性推导
&通过带容量约束的穷尽枚举排除所有第一因子秩型，无须补全第二、第三因子。
&取等给出 $M_tP^{-1}M_s=\delta_{ts}M_t$；张量乘积公式将三个因子耦合起来，与必需的可逆性矛盾。\\
验证
&检查限制证书与完整秩型枚举。
&原始计算路线检查继承证书及重新生成的 CNF/DRAT 文件对；Lean 路线检查整数分支证书与结构性推导，全部依赖经内核重新检查。\\
进一步结果
&另证明了 $R_{\mathbb F_3}(\langle2,3,3\rangle)=15$。
&饱和引理还约束零超额的秩 22 分解：至多一个第一因子可逆。所选位置的全部 22 个因子均非零时，秩型为 $(17,5,0)$ 或 $(18,3,1)$；含零因子的替代见第\ref{sec:outlook}节。\\
\end{longtable}

两条路线对共同的陪集几何采用了不同的处理：Wang 依据加强后的表检验第一因子容量，本文则强迫秩不等式取等，再导出跨因子恒等式。后者还给出第\ref{sec:outlook}节中的零超额秩型约束。

\paragraph{研究发展与开放材料。}
Wang 提供了包括代码和证书的可复现定理包。本报告还记录了结构性证明如何形成：秩 22 构造与形变、商核心、支撑与补全检验、纠错、回到完整张量，以及对饱和条件的识别。附录\ref{app:research-record}将这一发展整合为与证据相联的 Meta-Trace，并由主题研究档案支持。在结构性论证之外，这也构成对自主数学研究进行记录和研究的贡献。

对秩一张成搜索、对称性约化的秩型枚举或三元域矩形情形感兴趣的读者，可以在 Wang 的文章\cite{wang2026twentyone}中找到互补方法。结合本文的饱和论证与 Meta-Trace 阅读，该工作为理解有限域约束如何导出张量秩下界，提供了另一种视角。
